\chapter{二次型；初等几何.}

以现代面目出现的二次型理论，其历史很难追溯到 18 世纪下半叶之前，而且我们将会看到，它的发展主要是为了回应数论、分析和力学的需要。但这一理论的基本观念实际上从“欧几里得”几何的开端起就已经出现，并且构成了这种几何的框架。因此，若不至少以概述的方式谈谈自古以来“初等几何”的发展，便无法追溯它的历史。当然，我们只能勾勒少数几个一般观念的演变，读者不应指望在这里找到关于这个或那个具体定理的历史的精确资料，对此只需求助于专门的历史著作或教学著作即可。\(^{1}\)同样不言而喻的是，当我们在下文谈到同一定理用各种代数语言或几何语言表达时的各种可能解释时，我们完全无意宣称这些“翻译”在各个时代都像今天这样为人们所熟悉；恰恰相反，本篇笔记的主要目标正是要表明，数学家们是如何极其缓慢地逐渐意识到这些外表往往迥然不同的问题之间的联系；我们还必须表明，他们在这样做时，如何被引导着给古人遗留下来的一大堆几何定理注入某种连贯性，并最终试图精确地划定“几何”一词究竟应当理解为什么。

如果撇开巴比伦人关于二次方程求解公式的发现（{[}232{]}，第 183–189 页）不谈，那么就应当在几何外衣之下去寻找二次型理论诸主要概念的诞生。这些概念最初以距离的平方（在平面上或三维空间中）的形式出现，与之对应的“正交性”概念则借助直角引入，欧几里得把直角定义为平角的一半（《几何原本》第 I 卷，定义 10）；距离与直角这两个概念由毕达哥拉斯定理——欧几里得大厦的拱顶石——联系在一起。\(^{2}\)角的概念在希腊数学中似乎很早就被引入了（这无疑是得自巴比伦人，后者凭借长期的天文经验已经习惯于使用角）。众所周知，在古典时期，只有小于两直角的角才有定义（欧几里得的这个“定义”与他给出的直线和平面的定义一样，既含混又不可用）；方向的概念并未被提炼出来，尽管欧几里得使用了（既无公理亦无定义）直线把平面分成两个区域这一事实，而且在必要之处他还会仔细地区分这两个区域。\(^{3}\)在这个阶段，平面旋转群的观念只能以非常不完善的方式显露出来，即通过半线的非定向角的加法（它同样是欧几里得未加解释地引入的），而这种加法原则上只在和至多等于两直角时才有定义。\(^{4}\)至于三角学，它受到几何学家的轻视，被交给土地测量者和天文学家；正是后者（阿里斯塔克斯、喜帕恰斯，尤其是托勒密{[}255{]}）建立了直角三角形（平面的或球面的）边与角之间的基本关系，并编制了最初的数表（这些表给出半径为 r 的圆上被角 \(\theta < \pi\) 截出的弧所对的弦，换句话说就是数 \(2r \sin \frac{\theta}{2}\)；更便于使用的正弦的引入应归功于中世纪的印度数学家）；在计算这些表时，当时尚不为人知的弧加法公式被换成对托勒密定理的等价使用（它也许可以追溯到喜帕恰斯）

关于圆内接四边形的定理。还应当指出，欧几里得和海伦就任意平面三角形的边与角给出了与公式

\[
a ^ {2} = b ^ {2} + c ^ {2} - 2 b c \cos A
\]

等价的命题；但由于缺乏向量演算的观念——它要到 19 世纪才会出现——人们很难在其中看出与度量型相伴的双线性型概念的初次现身。

位移（或者说运动，这两个概念的区别在古代——甚至在很久之后——都不清楚）已为欧几里得所知；但出于我们无从知晓的原因，他似乎对使用位移怀有明显的反感（例如在“三角形全等的情形”中，人们得到的印象是，他之所以使用位移概念，只是因为他未能表述出一条适当的公理（{[}153 e{]}，第 I 卷，第 225–227 页及第 249 页））；然而，为了定义回转锥面和球面（《几何原本》第 XI 卷，定义 14 和 18），他求助于的恰恰是位移（绕一条轴的旋转）概念，阿基米德定义回转二次曲面时用的也是这一概念。但是，作用于整个空间的变换这一一般观念，在 18 世纪末之前几乎与数学思想无缘；⁵而在 17 世纪之前，人们既找不到运动合成概念的任何痕迹，更谈不上位移的合成。当然，这并不是说希腊人没有特别意识到图形的“正则性”和“对称性”——我们现在把这些概念与位移群的概念联系在一起：他们关于正多边形、尤其是正多面体——其全部数学中最出色的篇章之一——的理论，正好证明情况相反。⁶

最后，就我们所关心的领域而言，希腊数学最后一项本质贡献是圆锥曲线的理论（至于二次曲面，希腊人只知道某些回转二次曲面，而且除球面外，他们对这些曲面的研究都没有走得很远）。这里值得注意的是，尽管希腊人从未产生解析几何基本原理的想法（主要因为缺少一种便于使用的代数），他们在研究特殊“图形”时通常使用相对于平面上两条（甚至多于两条）轴的“纵标”（这些轴与图形紧密相连，这正是他们的方法与费马和笛卡尔的方法的根本差异之一，后者的轴是独立于所考虑的图形而固定的）。特别地，在倍立方问题中出现的第一批圆锥曲线（圆除外）的例子，是由方程 \(y^{2}=ax, y=bx^{2}, xy=c\) 给出的曲线（欧多克索斯的学生梅内克缪斯，公元前 4 世纪中叶）；\(^{7}\)而在研究与这些曲线有关的问题时，最常使用的是圆锥曲线的方程（通常相对于由一条直径以及曲线与它的一个接触点处的切线所组成的两条斜轴）（而“焦点”性质只起着非常不显眼的作用，这与仅仅上溯到 19 世纪的渊博传统会使我们相信的恰好相反）。在这门浩繁的理论中，我们在这里尤其应当记住共轭直径的概念（阿基米德已经知道它），以及今天用来定义一点的极线的那条性质，它由阿波罗尼奥斯{[}153 b{]}就该点在圆锥曲线之外的情形给出（因此对他而言，极线就是由从该点引出的各切线的切点连成的直线）；从我们的观点看，这是关于异于度量型的二次型的“正交性”的两个例子，但可以想见，这些概念与经典的垂线概念之间的联系，在那个时代是不可能被设想到的。

在笛卡尔和费马之前，几乎没有别的进展可以指出；但随着解析几何的开端，二次型的代数理论开始从它的几何外壳中脱胎而出：费马知道平面上的二次方程表示一条圆锥曲线（{[}109{]}，第 I 卷，第 100–102 页；法文译本，第 III 卷，第 84–101 页），并对二次曲面勾画了类似的想法（{[}109{]}，第 I 卷，第 111–117 页；法文译本，第 III 卷，第 102–108 页）。随着 18 世纪中二维和三维解析几何的发展，该理论的两个中心问题出现了（尤其是联系着圆锥曲线和二次曲面）：把二次型化为平方和，以及寻求它关于度量型的“轴”。对圆锥曲线来说，这两个问题太初等，引不起重要的代数进展；对任意多个变量的情形，第一个问题由拉格朗日于 1759 年解决，与多元函数的极大值有关（{[}191{]}，第 I 卷，第 3–20 页）。但这个问题几乎立刻被求轴的问题所掩盖，甚至早在秩的不变性被表述之前就是如此；\(^{8}\)至于惯性定律，它直到 1850 年前后才由雅可比（{[}171{]}，第 III 卷，第 593–598 页）发现，他以与今天相同的论证证明了它，而西尔维斯特（{[}304{]}，第 I 卷，第 378–381 页）则满足于把它当作近乎显然的事实加以陈述。\(^{9}\)

把二次曲面化为它的轴的问题，已经呈现出比圆锥曲线的相应问题大得多的代数困难；最先着手处理这个问题的欧拉，尚无力证明特征值的实性，他在一次缺乏说服力的论证尝试之后便承认了这一点（{[}108 a{]}，(1)，第 IX 卷，第 379–392 页）。\(^{10}\)这一点大约在 1800 年{[}140{]}才被正确建立，而关于任意 n 个变量的型的相应定理的证明，则必须等到柯西（{[}56 a{]}，(2)，第 IX 卷，第 174–195 页）。也是柯西，在大约同一时间，证明了给出特征值的特征方程在一切正交换轴之下不变（{[}56 a{]}，(2)，第 V 卷，第 252 页）；\(^{11}\)但对 n = 2 或 n = 3，这种不变性凭特征值由相应的圆锥曲线或二次曲面的轴给出的几何解释，在直观上是“显然的”。此外，在关于这个课题的研究过程中，特征值的初等对称函数也自然地出现过（带有种种几何解释，特别是与阿波罗尼奥斯关于共轭直径的定理相联系），尤其是判别式，它（在 n = 2 的情形，联系着二次方程理论而早已为人所知）在 n = 3 的情形由欧拉首次引入（{[}108 a{]}，(1)，第 IX 卷，第 382 页）；后者是在给二次曲面分类时遇到它的（他借此表达二次曲面没有无穷远点的条件），却未提及它关于正交换轴的不变性。但稍后，随着整系数二次型的算术理论的开端，拉格朗日（就 n = 2）注意到了判别式在线性但非正交的变量代换之下不变性的一个特例（{[}191{]}，第 III 卷，第 699 页），高斯则就 n = 3 建立了判别式对一切线性变换的“协变性”（{[}124 a{]}，第 I 卷，第 301–302 页）。\(^{12}\)一旦柯西和比内证明了行列式乘法的一般公式，把高斯的公式推广到任意多个变量的情形便是直接的了；正是这一点，在大约 1845 年，给了一般的不变量理论以最初的推动。

在希腊人那里取代位移理论的两个概念——运动的概念和图形“对称”的概念——之外，17 和 18 世纪又加上了第三个，即直角坐标轴变换的问题，它与位移理论实质上等价。欧拉为这个问题撰写了好几部著作，他最看重的是为换轴公式求得便于使用的参数表示。力学日后对他为 n = .3 的情形为此目的引入的那三个角作了什么用处，这是人所共知的（{[}108 a{]}，(1)，第 IX 卷，第 371–378 页）。但他并不止步于此：他在 1770 年考虑了任意 n 的正交变换的一般问题，指出这里通过引入 \(n(n - 1)/2\) 个角作为参数即可达到目的，最后，对 n = 3 和 n = 4，他给出了旋转的有理表示（分别是 4 个齐次参数、以及由一个关系式联系起来的 8 个齐次参数的函数），它们正是后来借助四元数理论得到的那些表示，而他并未指出其来历（{[}108 a{]}，(1)，第 VI 卷，第 287–315 页）。\(^{13}\)

另一方面，欧拉还指出了如何解析地表达对平面图形“对称”的寻求，正是为了这个目的，他被引向实质性地证明：一个平面位移是一个旋转，或是一个平移，或是一个平移后面接一个对称（{[}108 a{]}，(1)，第 IX 卷，第 197–199 页）。力学的兴起此时还导致对位移的一般研究：但起初涉及的只是与连续运动相切的“无穷小”位移：在托里拆利、罗贝瓦尔和笛卡尔关于运动合成以及平面运动的瞬时旋转中心的研究中，出现的显然只是这种位移（参见第 177 页）。后者由约翰·伯努利给出了一般的定义；达朗贝尔于 1749 年、欧拉于次年推广了这个概念，证明了对使一点固定不动的运动存在瞬时旋转轴。有限位移的相应定理直到 1775 年才由欧拉{[}108 b{]}陈述出来，在他于此同时发现旋转的行列式等于 1 的那篇论文中；次年，他证明了平面相似变换存在不动点（{[}108 a{]}，(1)，第 XXVI 卷，第 276–285 页）。但要最终得到关于有限位移和无穷小位移的一个连贯理论，还须期待沙勒始于 1830 年的工作{[}60 a{]}。

这样我们就来到了可称之为几何学黄金时代的时期，它大致处于蒙日《画法几何》（1795 年）{[}225{]}的出版与 F. 克莱因的“埃尔朗根纲领”（1872 年）（{[}182{]}，第 I 卷，第 460–497 页）这两个日期之间。这一几何学的骤然复兴带给我们的本质进展如下：

\begin{enumerate}
\def\labelenumi{\Alph{enumi})}
\item
  无穷远元素（点、直线或平面）的概念由德萨格在 17 世纪引入{[}84{]}，但在 18 世纪里，它除了被滥用为一种说法之外几乎没有得到表达；彭赛列{[}252{]}恢复了它的地位并系统地使用它，从而把射影空间弄成了一切几何现象的一般框架。
\item
  与此同时，随着蒙日、尤其是彭赛列的工作，向复射影几何的过渡得以完成。在 18 世纪中只被零星使用的虚点概念，在这里（与无穷远点的概念同时）得到利用，以给出独立于实仿射几何的“图形情形”的命题。如果说最初为支持这些创新而提出的辩护还很笨拙（尤其是出自“综合”几何学派成员之手的辩护，在这个学派中，坐标的使用最终被视为一件不体面的事），那么不容错过的是，在那里可以辨认出，以蒙日的“偶然关系原理”或彭赛列的“连续性原理”的名义出现的，正是现代代数几何“特殊化”思想的最初萌芽。\(^{14}\)
\end{enumerate}

从这些概念得出的最早结果之一，是注意到在复射影空间中，所有非退化的圆锥曲线（相应地，二次曲面）都属于同一类；这引导彭赛列发现了“迷向”元素：“平面上任意放置的各个圆”，他说，“因此并不是彼此完全独立的，不像初看时所可能认为的那样，它们在理想意义下于无穷远处有两个公共的虚点”（{[}252{]}，第 I 卷，第 48 页）。后来，他又以同样的方式引入了“脐形曲线”，即所有球面公共的无穷远虚圆锥曲线（{[}252{]}，第 I 卷，第 370 页）；即使他没有特别谈及球面的迷向母线，他至少明确强调了所有二次曲面都存在实的或虚的直母线（同前引文，第 371 页）；\(^{15}\)这些概念，他的后继者们（尤其是

普吕克和沙勒）对它们的使用比他本人还要多，特别是在圆锥曲线和二次曲面的“焦点”性质的研究中。

\begin{enumerate}
\def\labelenumi{\Alph{enumi})}
\setcounter{enumi}{2}
\item
  点变换和变换合成的概念也同样得到了一般的表述，并被系统地引进作为证明手段。除了位移和投影之外，直到那时人们只知道一些特殊的变换：拉伊尔和牛顿使用过的 \(x' = a / x\)、\(y' = y / x\) 型的某些平面射影变换，克莱罗和欧拉的“仿射变换”\(x' = ax\)、\(y' = by\)（{[}108 a{]}，(1)，第 IX 卷，第 XVIII 章），最后还有再次由牛顿以及麦克劳林和布雷肯里奇给出的若干特殊二次变换。蒙日在其《画法几何》中展示了平面投影在三维几何中能得到多么充分的利用。在彭赛列那里，一种已被用到泛滥的系统化证明程序，就是通过投影把圆锥曲线的性质化归为圆的性质（笛卡尔和帕斯卡已经偶尔使用过这种方法）；而为了从二次曲面过渡到球面，他发明了空间的射影变换的第一个例子，即“透射”（{[}252{]}，第 I 卷，第 357 页）；最后，也是他引入了曲线到自身的双有理变换的最初例子。1827 年，默比乌斯（{[}223{]}，第 I 卷，第 217 页）（以及 1830 年沙勒的独立工作（{[}60 b{]}，第 695 页））定义了最一般的线性射影变换；与此同时出现了反演以及其他类型的二次变换，对它们的研究将开创双有理变换的理论，这一理论将在 19 世纪下半叶发展起来。
\item
  对偶的概念完全公开登场，并且被发现自觉地与双线性型理论联系在一起。关于圆锥曲线的极点与极线的理论，自阿波罗尼奥斯以来只在德萨格和拉伊尔那里取得过一些进展，蒙日把它推广到了二次曲面，他和他的学生们都看到了借此把已知定理变换成新结果的可能性。\(^{16}\)但是，把这些评述建构成“互反配极”变换理论中的一般方法、并把它变成一种特别有效的发现工具，这一功劳仍然要归于彭赛列。稍后，特别是在热尔岗、普吕克、默比乌斯和沙勒那里，对偶的一般概念摆脱了它与二次型的联系（这种联系在彭赛列那里还过于紧密）。特别地，默比乌斯在考察三维空间中对偶（由一个双线性型定义）的各种可能性时，于 1833 年发现了关于交错双线性型的对偶（{[}223{]}，第 I 卷，第 489–515 页），\(^{17}\)它在 19 世纪尤其以“线性线丛”理论的形式得到研究，并联系着凯莱、格拉斯曼和普吕克约在 1860 年引入的“线几何学”和“普吕克坐标”而发展起来。
\item
  从射影几何的开端起，对经典几何的各种性质及其与射影空间的关系的深入研究，很快导致把它们划分为“射影性质”和“度量性质”；把这一区分视为日后成为现代结构概念的那种东西在当时最清晰的表现之一，无疑并不算夸张。但是最先引入这一区分及其术语的彭赛列，已经意识到联系着这两类性质的东西；在《论著》中处理关于角的问题时，这些性质“似乎不属于我们称之为射影的那类性质 \ldots，然而它们毕竟以如此简单的方式”，他说，“得自构成{[}本书{]}基础的那些原理 \ldots，以致我不认为有任何别的几何理论能够以既更直接又更简单的方式导出它们。如果考虑到图形的射影性质必然是图形所能具有的性质中最一般的性质，这就不会令人惊讶；这样一来，它们必定像简单的推论一样，把该领域所有其他的性质或特殊关系都包括进来”（{[}252{]}，第 I 卷，第 248 页）。说实话，在这番宣言之后，看到他以非常迂回的方式处理关于角的问题——把角与圆锥曲线的焦点性质联系起来，而不是直接引入圆点——未免有点令人惊讶；而事实上，直到 30 年后，拉盖尔（当时还是综合理工学校的学生）才借助这些直线与同起点的各迷向直线的交比给出直角的表达式（{[}192{]}，第 II 卷，第 13 页）。最后，随着凯莱（{[}58{]}，第 II 卷，第 561–592 页）的工作，平面图形的“度量”性质无非是该图形添上圆点之后的“射影”性质这一基本思想得到了明确的表达——这是通向“埃尔朗根纲领”的一个决定性中途站。
\item
  将在 1830 年前后问世的双曲型非欧几何，起初多少置身于我们正勾勒其主要脉络的那场运动之外。这门新几何产生于本质上属于逻辑层次、涉及经典几何基础的关切，其发明者\(^{18}\)以与欧几里得几何相同的公理化“综合”形式表述它，而且它与射影几何毫无联系（在这门几何中唯一平行的概念消失了，因此按照经典模式引入射影几何甚至似乎先验地就被排除）；无疑正是由于这个缘故，在很长时间里，它几乎引不起法国、德国或英国的射影几何学派的任何兴趣。于是，当凯莱在前引的那篇基本论文（{[}58{]}，第 II 卷，第 561–
\end{enumerate}

\begin{enumerate}
\def\labelenumi{\arabic{enumi})}
\setcounter{enumi}{591}
\tightlist
\item
  592 页）中想到用一条任意圆锥曲线（他称之为“绝对形”）来代替圆点（被看作“相切退化”的圆锥曲线）时，他根本没有想到把这个想法与罗巴切夫斯基–鲍耶的几何联系起来，尽管他指出了他的构想如何导出两点间“距离”的新表达式，也提到了它与球面几何的联系。局势大约在 1870 年发生变化，此时非欧几何随着罗巴切夫斯基著作的传播、高斯著作的出版以及黎曼的就职演讲而来到数学新闻的前沿。沿着黎曼开拓的道路，贝尔特拉米在不知道凯莱工作的情况下，于 1868 年重新找到了后者给出的距离表达式，但处在相当不同的语境中：他把一个圆的内部看作一张常曲率曲面的像，其中测地线由直线来表示{[}18 a{]}；两年后，克莱因（独立于贝尔特拉米）对这些不同的观点作了综合，并通过发现椭圆型非欧空间（{[}182{]}，第 I 卷，第 254–305 页）而使之完备。\(^{19}\)
\end{enumerate}

\begin{enumerate}
\def\labelenumi{\Alph{enumi})}
\setcounter{enumi}{6}
\tightlist
\item
  在我们此处考察的时期的后半段，还插入了一个批判性反思的阶段，在此期间，“综合”几何的支持者们不满足于把坐标从证明中驱逐出去，甚至打算在几何公理中也弃用实数。这一学派的主要代表是冯·施陶特，他基本上成功地完成了这一杰作{[}325{]}，它在其时代、乃至进入 20 世纪很久之后都备受推崇；虽然今天人们对这种层次的思想不再给予同样的重视——其富有成效的应用前景被证明相当有限——但仍应承认，冯·施陶特及其弟子们的努力有助于澄清实的或复的“标量”在经典几何中的作用，并恰恰以此引入了在任意基域上的几何这一现代观念。
\end{enumerate}

大约在 1860 年，“综合”几何达到了它的顶峰，但其统治的终结也在迅速临近。在整个 18 世纪一直笨重而粗陋的解析几何，在拉梅、博比利耶、柯西、普吕克和默比乌斯手中，终于获得了使它得以与对手平起平坐的优雅和简洁。尤其是，大约 1850 年以后，群和不变量的观念终于获得了精确的表述，一点点地占领了舞台，人们看到，经典几何的定理无非是相似群的不变量或共变式之间的恒等关系的表达，\(^{20}\)正像射影几何的定理表达射影群的共变式之间的恒等关系（即“合冲”）一样。正是 F. 克莱因在著名的“埃尔朗根纲领”（{[}182{]}，第 I 卷，第 460–497 页）中权威地陈述了这一论点，他在其中主张抛弃“综合”倾向与“解析”倾向之间毫无结果的争论；他说，如果说针对后者让任意的坐标系扮演特权角色所提出的指责，“就先前使用坐标方法时的缺陷而言，在大多数情况下是正当的，那么当问题在于对这种方法的合理运用时，它就站不住脚了……空间直观的领域并没有被解析方法所禁止……”，他还强调，“人们不应低估一个合用的格式给后续研究带来的好处，因为它可以说走在思想的前面”（同前引文，第 488–490 页）。

于是人们最终得到了按照定理所从出的群对“几何”定理所作的一种合理的“结构性”分类：射影几何对应线性群，度量问题对应正交群，“线性线丛”几何对应辛群。但在这道无情的强光之下，经典几何——代数几何和微分几何除外，\(^{21}\)它们从此被确立为自治的学科——骤然熄灭，失去了全部光彩。基于变换之使用的诸方法的推广，早已使新定理的产出变得有些机械：沙勒 1837 年在其《历史概览》中说：“今天，人人都可以挺身而出，任取一条已知的真理，把它交付给各种一般的变换原理；他将得出别的真理，不同的或更一般的；而这些又可接受类似的操作；于是，从第一条真理推导出来的新真理的数目几乎可以倍增到无穷 \ldots{} 因此，在科学的现状下，谁愿意，谁就能够加以推广并在几何中创造；为这座大厦添上一块石头，天才已不再是不可或缺的”（{[}60 b{]}，第 268–269 页）。但随着不变量理论的进展，局势变得清楚多了：它终于（至少对“经典”群而言）表述出了一些一般方法，原则上能以纯粹自动的方式写出一切代数共变式及其全部“合冲”；这是一场胜利，它同时宣告了经典不变量理论本身以及“初等”几何作为研究领域的死亡，\(^{22}\)后者在实践中已变成它的一本简单词典。无疑，在可以随意大量炮制出来的无穷多条“定理”当中，没有什么办法先验地预见哪些定理用恰当的几何语言陈述出来会具有可与经典结果相比拟的简洁与优美，那里还留有一块狭窄的领域，众多业余爱好者仍乐于在其中耕耘（三角形、四面体的几何，低次代数曲线与曲面的几何，等等）。但对职业数学家来说，矿藏已经采竭，因为那里不再有任何可能波及数学其他部分的结构性问题；群论与不变量理论的这一篇章，在另行通知之前，可以视为已经终结。\(^{23}\)

因此，在埃尔朗根纲领之后，欧氏几何与非欧几何从纯代数的观点看，变成了或多少便利的简单语言，用以表达双线性型理论的结果，而后者的进展与不变量理论的进展齐头并进。\(^{24}\)凡涉及双线性型的秩的概念以及这些型与线性变换之间的关系的一切内容，都由弗罗贝尼乌斯的著作（{[}119{]}，第 I 卷，第 343–405 页）得到了最终的澄清。也是多亏弗罗贝尼乌斯，我们才有了一个交错型在自由 Z-模上的标准表达式（{[}119{]}，第 I 卷，第 482–544 页）；不过，斜对称行列式早在本世纪初已随普法夫出现，联系着把微分形式化为规范形的问题；雅可比在 1827 年重新研究了这个问题（{[}171{]}，第 IV 卷，第 17–29 页），他知道奇数阶的斜对称行列式为零，而且正是他构造出普法夫式的表达式并证明它是偶数阶斜对称行列式的一个因子；但他没有意识到后者是普法夫式的平方，这一点直到 1849 年才由凯莱建立（{[}58{]}，第 I 卷，第 410–413 页）。与一个二次型相伴的对称双线性型的概念是“配极化”过程的最简单情形，而配极化是不变量理论的基本工具之一。以“标量积”的名义，这一概念将享有巨大的昌盛：先是在“向量演算”的普及者那里，然后从 20 世纪起，得益于希尔伯特空间理论带给它的未曾预料到的推广（见第 212 页）。也正是后一种理论将使算子的伴随这一概念显露出来（在此之前，它几乎只在线性微分方程理论中露过面，而在张量演算中，它则表现为度量张量的指挥下协变指标与逆变指标的轮舞）；最终赋予埃尔米特型概念以其全部个性的也正是这一理论，该概念最初由埃尔米特于 1853 年在其算术研究中引入（{[}159{]}，第 I 卷，第 237 页），但直到大约 1925 年以及复希尔伯特空间应用于量子理论之前，它一直处于数学大潮的边缘。

正交群和相似群——自 19 世纪中叶起就被清楚地设想并照此处理，而且已经成为二次型理论的核心——以及其他“经典”群（线性群、辛群和酉群）的研究，则占据了越来越重要的地位。这些群在李群理论与微分几何中、以及在二次型算术理论（例如见{[}285{]}和{[}100{]}）中所起的重要作用，我们在此只能提上一句；\(^{25}\)正是由于这一情况，也由于对偶概念向最多样问题的推广，在现代数学理论中才几乎没有哪一门不以这种或那种形式牵涉到双线性型。无论如何应当指出，正是（三维）旋转群的研究引导哈密顿发现了四元数{[}145 a{]}；这一发现由 W. 克利福德加以推广，他于 1876 年引入了以他的名字命名的那些代数，并证明它们是若干四元数代数、或四元数代数与一个二次扩张的张量积（{[}65{]}，第 266–276 页）。这些代数四年后被利普希茨重新发现{[}205 b{]}，他用它们给出了 n 个变量的正交变换的参数表示（推广了凯莱借助四元数理论就 n = 3（{[}58{]}，第 I 卷，第 123–126 页）和 n = 4（第 II 卷，第 202–215 页）得到的那些表示），它们以及由之派生的“旋量”概念（{[}52 b{]}和{[}62 a{]}）也因其在量子理论中的应用而在现代享有巨大的流行。

最后，我们还须谈一谈思想的演变，它导致半双线性型理论中关于标量环的一切限制几乎全部被抛弃——这是整个现代代数中共同的倾向，但它在这里也许比在别处显现得更早。我们已经指出过复数域上的几何的富有成效的引入（此外，在整个 19 世纪，这种几何与实几何之间一直存在一种持久的、有时是危险的混淆）；这方面的清晰特别来自 19 世纪末关于几何基础的公理化研究{[}163 c{]}。在这类研究中，希尔伯特及其追随者在考察各条公理之间的关系时，被引去构造适当的反例，其中的“基域”（交换的或非交换的）具有或多或少的病态性质，他们由此使数学家们习惯于全新类型的“几何”。从分析的观点看，伽罗瓦已经考虑过系数和变量取值于有限素域的线性变换（{[}123{]}，第 145 页）；在发展这些思想时，若尔当{[}174 a{]}很自然地被引导去考虑这些域上的经典群，这些群的介入显现在各种各样的数学领域中。迪克森在大约 1900 年把若尔当的研究推广到一切有限域，而更近一些，人们已经看到，若尔当–迪克森理论的一大部分可以推广到绝对任意的“基域”上；这本质上应归功于迷向向量的一般性质和维特定理，它们在经典情形是平凡的，直到 1936 年才就任意基域建立起来{[}337 a{]}。\(^{26}\)

但是，在把半双线性型的研究推向不断增大的“抽象”的同时，人们发现，把二维或三维情形中出自经典几何的术语原样保留下来，并把它推广到 n 维情形乃至无穷维空间，是极富启发性的。经典几何就这样卸下了它作为一门自治的、有生命的科学的角色，被重塑为当代数学的一种普遍语言，其灵活性和有用性都是无与伦比的。

\chapter{拓扑空间.}

极限和连续的概念可以追溯到古代；若不从这个观点出发系统地研究不仅数学家们、而且还有希腊哲学家们特别是亚里士多德，若不追随这些思想在文艺复兴时期数学以及微分积分学开端中的演变，便不可能写出一部完整的历史。这样一项研究做起来固然非常有意思，但它会远远超出本篇笔记的框架。

应当把拓扑学——如同现代数学的许多其他分支一样——的创立者视为黎曼：事实上正是他第一个力图提炼出拓扑空间的概念，构想出关于这些空间的一门自治理论的思想，定义了在拓扑学日后发展中起最大作用的不变量（“贝蒂数”），并给出了它对分析的首批应用（阿贝尔积分的周期）。但 19 世纪上半叶的思想运动并非没有在不止一个方面为黎曼做好准备。的确，把数学安放在坚实基础之上的愿望——它是整个 19 世纪直至我们今天如此众多重要研究的起因——已经导致了收敛级数的概念和趋于一个极限的数列的概念（柯西、阿贝尔）以及连续函数的概念（波尔查诺、柯西）的正确定义。另一方面，复数——或者像当时人们所说的那样，“虚数”（在 18 世纪有时还被冠以“不可能的数”之名）——的几何表示（用平面上的点），这种归功于阿尔冈和高斯（见第 161 页）的表示法，已为大多数数学家所熟悉：它所构成的进步，与我们今天在希尔伯特空间研究中采用几何语言属于同一层次，并且以胚芽的形态包含着对每一个可连续变化的对象作几何表示的可能性；高斯本来就因其关于几何基础、非欧几何、曲面的研究而很自然地被引向这类观念，他看来已经预见到了这种可能性，因为他在（独立于阿尔冈和法国数学家们）定义虚数的几何表示时使用了“二重广延的大小”一语（{[}124 a{]}，第 II 卷，第 101–103 页及第 175–178 页）。

一方面是他关于代数函数及其积分的研究，另一方面是他关于几何基础的思考（在很大程度上得自他对高斯著作的研读），正是这两者引导黎曼表述出一个恰好就是现代拓扑学研究纲领的纲领，并使这一纲领的实现有了开端。例如，他在其《阿贝尔函数论》（{[}259 a{]}，第 91 页）中是这样说的：

“在研究通过积分恰当微分而得到的函数时，位置分析的几条定理几乎是不可缺少的。这个名称曾为莱布尼茨所用，尽管含义也许有些不同，我们要把它理解为连续量理论中研究这些量的那一部分：不是把这些量看作独立于其位置、可用彼此量度的东西来研究，而是抽掉一切度量的观念，只研究它们的位置关系和包含关系。我保留日后以一种完全独立于一切度量的方式处理这个对象的权利 \ldots{}”。

而在他著名的就职演讲《论作为几何基础的假设》（{[}259 a{]}，第 272 页）中：

“\ldots 多次广延的一般大小概念，\(^{1}\)它把空间大小作为特殊情形包含在内，至今完全没有得到研究 \ldots{}（第 272 页）”。

“\ldots 大小的概念预设了一个可取各种不同规定的元素。根据能否从一个规定连续地过渡到另一个规定，这些规定构成一个连续流形（它们将被称为该流形的点）或一个离散流形（第 273 页）”。

“\ldots 度量在于把要比较的各个量相互叠置；因此，要进行度量，就必须有一种把一个量移置到另一个量上的手段。在缺乏这种手段时，只有当一个量构成另一个量的一部分，才可能比较这两个量\ldots{}就此所能作的研究，构成了大小理论中独立于度量的那一部分，在这里，量既不被看作独立于其位置而存在，也不被看作可用一个度量单位来表达，而被看作一个流形的各个部分。这样的研究在数学的许多部分中已成为必需，对解析多值函数的理论来说尤其如此 \ldots{}（第 274 页）”。

“\ldots 在一个给定的流形中，位置的确定，凡属可能，都归结为有限多个数值规定。诚然，存在这样一些流形，其中位置的确定所需要的不是有限多个规定，而是一个无穷序列或一个连续的量规定之流形。这样的流形例如由一个函数在一个给定区域中的各种可能规定、由空间中一个图形的位置等等构成”（第 276 页）。

在这最后一句里，可以注意到函数空间研究的最初想法；此外，早在黎曼的博士论文中，同样的想法就已经得到表达：他在谈到以狄利克雷原理之名著称的最小值问题时说，“这些函数的全体构成一个连通的、自身封闭的区域”（{[}259 a{]}，第 30 页），这虽然形式尚不完善，却已经是希尔伯特日后给出的狄利克雷原理证明、乃至函数空间对变分法的大多数应用的胚胎。

如前所述，黎曼为这一宏伟纲领的实现开了头：他定义了“贝蒂数”，先是一个曲面的贝蒂数（{[}259 a{]}，第 92–93 页），然后（同前引文，第 479–482 页；另参见{[}259 c{]}）是任意维数的一个流形的贝蒂数，并把这一定义应用于积分理论；这些结果开创了代数拓扑学，这门数学分支自 20 世纪初以来发展不断加速，其历史不应在此追述。

至于拓扑空间的一般理论，即黎曼所预见的那种理论，它要得到发展，就必须先对实数理论、数集理论、直线上的点集理论以及平面和空间中的点集理论，作比黎曼时代更为系统的研究：这一研究一方面联系着关于无理数本性的探讨（波尔查诺的偏于哲学，戴德金的则本质上是数学的），另一方面联系着实变量函数论的进展（黎曼本人以他的积分定义和他的三角级数理论为它作出了重要贡献，杜布瓦-雷蒙、迪尼、魏尔斯特拉斯等人的工作也以它为对象）；这是 19 世纪下半叶的工作，而其中最关键的是康托尔，他第一个定义了（先是在直线上，然后在 n 维欧几里得空间中）聚点、闭集和完全集的概念，并得到了关于这些集在直线上的结构的本质结果（参见第 155 页）：对此不仅应查阅康托尔的著作{[}47{]}，还应查阅他与戴德金的那批非常有趣的通信{[}48{]}，在那里还可以清楚地看到把维数视为拓扑不变量的思想。

理论的后续进展，例如在舍恩弗利斯的书{[}275 a{]}中以半历史、半系统的形式作了陈述：其中远远最为重要的是博雷尔–勒贝格定理（最先由博雷尔就直线上的一个闭区间和覆盖它的可数开区间族证明）。

康托尔的思想起初遇到了相当激烈的反对（参见第 28 页）。但至少他关于直线上和平面上的点集的理论，很快就被法国和德国的函数论学派（若尔当、庞加莱、克莱因、米塔-列夫勒，随后是阿达马、博雷尔、贝尔、勒贝格等）所使用和广泛传播：特别是博雷尔丛书的开头几卷，每一卷都含有这一理论的一个初等陈述（例如见{[}32 a{]}）。随着这些思想的传播，人们开始以各种方式思考把它们应用于不再是点集、而是曲线之集或函数之集的可能性：例如早在 1884 年，这一想法就已见诸阿斯科利一篇论文的标题“论曲线族的极限曲线”{[}10{]}，并且它还表达在阿达马 1896 年向苏黎世数学家大会所作的一次通报中{[}141{]}；它也与沃尔泰拉 1887 年引入“线函数”以及创立“泛函演算”（即变元为函数的函数论）紧密相联（关于这一点，可参看沃尔泰拉关于泛函分析的著作{[}322 b{]}）。另一方面，在那篇著名的论文（{[}163 a{]}，第 III 卷，第 10–37 页）中，希尔伯特在这一点上重拾黎曼的思想，证明了狄利克雷原理中最小值的存在，并开创了变分法的“直接方法”，在其中可以清楚地看到考虑使波尔查诺–魏尔斯特拉斯原理成立的函数集（也就是说，每个序列都含有一个收敛的子序列）有何意义；这样的集合不久确实发挥了重要作用，不仅在变分法中，而且在实变量函数论（阿斯科利、阿尔泽拉）和复变量函数论（维塔利、卡拉泰奥多里、蒙泰尔）中都是如此。最后，函数方程的研究，尤其是弗雷德霍姆对以他的名字命名的那类方程的求解{[}116{]}，使人们熟悉了把函数看作变元、把函数之集看作点集之类似的考虑，对这样的集合使用几何语言，与对 n 维欧几里得空间的点使用几何语言同样自然（这种空间同样也逃出了“直观”，并因此而在很长时间里始终是许多数学家不信任的对象）。特别地，希尔伯特关于积分方程的那些值得纪念的工作{[}163 b{]}，以埃哈德·施密特{[}274 b{]}对希尔伯特空间的定义和几何研究而告终，它与欧几里得几何完全类似（见第 212 页）。

与此同时，公理化理论的概念已变得愈来愈重要，这首先要归功于关于几何基础的大量著作，其中希尔伯特的{[}163 c{]}产生了尤其决定性的影响；正是在这项工作期间，希尔伯特早在 1902 年（{[}163 c{]}，第 180 页）就被引导去精确陈述黎曼意义上的“二重广延流形”的第一个公理化定义，照他的说法，这个定义构成“对位置分析作严格公理化处理的基础”，而且已经使用了邻域（其含义受到希尔伯特当时所限定的问题之要求的限制）。

把点集的性质与函数集的性质中共同的东西提炼出来（而不引入“距离”的概念）的最初尝试，是由弗雷歇{[}115 a{]}和 F. 里斯{[}260 b{]}作出的；但前者从可数极限的概念出发，未能为不可度量化的空间构造出一个合用而富有成果的公理体系；至少他认识到了波尔查诺–魏尔斯特拉斯原理与博雷尔–勒贝格定理之间的关系；正是在这一关联中他引入了“紧”这个词，尽管其含义与今天赋予它的含义略有不同。至于 F. 里斯，他从聚点（或者更确切地说，从与之等价的“导集”）的概念出发，他的理论仍然是不完全的，况且始终只是一个概要。

随着豪斯多夫（{[}152 a{]}，第 7-8-9 章），今天所理解的一般拓扑学开始了。他重新拾起邻域的概念，并且懂得从希尔伯特关于平面上邻域的公理中挑选出那些能够同时赋予他的理论以一切想要的精确性和一般性的公理。他推演这些公理的推论的那一章，至今仍是公理化理论的一个典范：抽象，然而从一开始就适用于应用。很自然，这是一般拓扑学后续研究的出发点，尤其是莫斯科学派的工作，其大部分指向度量化问题（参见第 166 页）：我们在这里特别应当从中记住亚历山德罗夫和乌雷松对紧空间的定义（以“双紧空间”之名），然后是吉洪诺夫{[}313{]}关于紧空间的乘积的紧性的证明。最后是 H. 嘉当{[}53{]}引入的滤子，它为各种各样的应用提供了一个非常宝贵的工具（在其中它是“摩尔–史密斯收敛”概念{[}227{]}的一个有利的替代品），并且借助关于超滤子的定理，最终实现了理论的澄清与简化。

\chapter{一致空间.}

涉及一致空间的主要概念和命题是从实变量理论中一点点脱胎出来的，直到晚近才成为系统研究的对象。柯西为了严格地建立级数理论（参见第 153 页），把一条他似乎视为显然的原则当作出发点：序列 \((a_{n})\) 收敛的一个必要充分条件是，\(|a_{n+p} - a_{n}|\) 只要当 \(n\) 充分大就可任意小（例如见（{[}56 a{]}，(2)，第 VII 卷，第 267 页））。他与波尔查诺{[}27 c{]}无疑都属于最早明确陈述这一原则并认识到其重要性的人：因此，满足上述条件的实数序列被称为“柯西序列”，而且引申开来，度量空间中这样的点列 \((x_{n})\) 也获此名：\(x_{n+p}\) 与 \(x_{n}\) 之间的距离只要当 \(n\) 充分大就可任意小；最终，一致空间中柯西序列的推广便得到了“柯西滤子”之名。

后来，当实数的直观概念不再令人满意、人们为了给分析以坚实的基础而设法从有理数出发定义实数时，在 19 世纪下半叶提出的各种定义中，恰恰是柯西原则提供了最富成效的一种；这就是康托尔的定义（{[}47{]}，第 93–96 页）（海涅{[}154 b{]}等人也从康托尔的思想出发发展了它，梅雷则独立地作了发展），按照这个定义，让每个有理数的柯西序列（用康托尔的术语说，即“基本序列”）对应于一个实数；两个有理数柯西序列 \((a_{n})\) 与 \((b_{n})\) 对应于同一个实数，当且仅当 \(|a_{n} - b_{n}|\) 趋于零。这里的本质思想是，从某种观点看，有理数集 \(\mathbf{Q}\) 是“不完备的”，而实数集是由 \(\mathbf{Q}\) 经“完备化”而导出的那个“完备”集。

另一方面，海涅在主要受魏尔斯特拉斯和康托尔的思想启发的工作中，第一个定义了一元或多元实变量数值函数的一致连续性{[}154 a{]}，并证明了在 R 的一个闭有界区间上连续的每个数值函数都在其上一致连续{[}154 b{]}：这就是“海涅定理”。这一结果联系着 R 中闭有界区间的紧性（“博雷尔–勒贝格定理”，参见第 141 页），而且海涅对其定理所给的证明略作修改也可用来证明博雷尔–勒贝格定理（在好几位作者看来，这就足以成为把这个定理命名为“海涅–博雷尔定理”的理由）。

这些思想向更一般空间的推广，是在人们先研究一些特殊情形、然后一般地研究度量空间时完成的；在度量空间中给定了距离（点对的数值函数，满足某些公理），它同时定义了一个拓扑和一种一致结构。弗雷歇第一个陈述了这些空间的一般定义，认识到了柯西原则的重要性{[}115 a{]}，并且还就一切度量空间引入了预紧（或“全有界”{[}115 a 与 b{]}）空间的概念。豪斯多夫在其《集合论》（Mengenlehre）{[}152 a 与 b{]}中发展了度量空间理论的许多内容，他特别认识到，可以把上面考察过的康托尔构造应用于这些空间，从而对每个非“完备”的度量空间（也就是说柯西原则不成立的度量空间）导出一个“完备”的度量空间。

度量空间是一种特殊类型的“一致空间”；一致空间直到最近才由 A. 韦伊{[}330 b{]}以一般的方式加以定义。在此之前，人们只知道在处理度量空间时如何使用涉及“一致结构”的概念和结果：这就解释了在现代拓扑学的许多工作中，度量空间或可度量化空间（特别是紧的可度量化空间）为何在一些距离并无真正用处的问题里扮演着重要角色。一旦陈述了一致空间的定义，把例如豪斯多夫所陈述的度量空间理论的几乎全部内容推广到这些空间，便没有任何困难（尤其是当滤子的概念也可供使用时）（而且可以用同样的方式，把亚历山德罗夫–霍普夫的《拓扑学》{[}4{]}中就紧度量空间所陈述的结果，例如推广到一切紧空间）。特别地，一致空间的完备化定理无非是康托尔关于实数的构造的移植，不包含任何本质的修改。

\chapter{实数.}

对大小的每一次量度都蕴含着实数的一个弥散的观念。从数学的观点看，实数理论的起源必须追溯到巴比伦科学中一种记数系统逐步形成的过程，这种系统（原则上）能够把值标写得任意接近每一个实数{[}232{]}。掌握了这样一种系统，以及由此必然随之而来的对数值计算的信心，事实上不可避免地以实数的一种“朴素”概念而告终；这种概念与今天在初等教学中、或在物理学家和工程师那里（与十进记数法相联系）重新见到的概念几乎没有差别；这一概念不容许被精确地定义，但可以这样来表述：一个数被认为是通过获得近似值并把它们纳入计算的可能性来定义的：此外，这必然蕴含着由经验给出的、自然不容许无限逼近的大小的量度，与诸如 \(\sqrt{2}\) 之类的“数”（假定存在无限逼近这个数的算法）之间的某种程度的混淆。

因此，类似的“实用主义”观点会在一切以计算技巧取代对严格性的追求和理论关切的数学学派中重新出现。相反，正是后面这些东西支配着希腊数学：我们得感谢它给出了第一个严格而连贯的量的比的理论，也就是说，本质上是实数的理论；它是一系列关于比、特别是关于不可公度比的发现的最终结果，这些发现在希腊思想史上的重要性怎么强调也不为过，但由于缺乏确切的文本，我们只能辨认出其主要轮廓。希腊数学在其开端与关于比例、相似和比的思辨——部分是科学的，部分是哲学的和神秘的——不可分割地联系在一起，特别是与“简单比”（可用分子分母都很小的分数表示的比）的思辨；而宣称用整数和整数之比解释一切，正是毕达哥拉斯学派的诸种特色倾向之一。但恰恰是毕达哥拉斯学派发现了正方形的边与其对角线的不可公度性（\(\sqrt{2}\) 的无理性）：这无疑是数学中不可能性证明的第一个例子；仅仅陈述这个问题这一事实，就已蕴含着比与它的近似值之间的明确区分，并足以指出把希腊数学家与其前辈分隔开来的巨大鸿沟。\(^{1}\)

关于伴随这一重要发现及其后继的思想运动，我们知之甚少。\(^{2}\)我们只打算概括地指出量的比理论赖以建立的主要思想；这个理论由大数学家欧多克索斯（柏拉图的同时代人和朋友）整理定型，被古典希腊数学最终采纳，并通过欧几里得的《几何原本》{[}107{]}为我们所知，它在其中（《几何原本》第 V 卷）得到了大师式的阐述：

\begin{enumerate}
\def\labelenumi{\arabic{enumi})}
\item
  数这个词和数的观念被严格保留给 \textgreater{} 1 的整数（1 是单子，而不是严格意义上的数），这不仅把我们的无理数排除在外，甚至也把我们称为有理数的那些数排除在外，因为对古典时代的希腊数学家来说，它们只是数之比。这里远不止一个单纯的术语问题：“数”这个词对希腊人（以及对直到近代的现代人）来说，是与具有双重合成法则（加法与乘法）的系统的观念联系在一起的：古典希腊数学家把整数之比设想为算子，定义在整数集或其一部分之上（p 与 q 之比是这样的算子：当 N 是 q 的倍数时，它让 N 对应于整数 \(p.(N/q)\)），这些算子构成一个乘法群，但不是具有双重合成法则的系统。在这一点上，希腊数学家自觉地与“计算术家”即职业计算者划清界限，后者像他们的埃及和巴比伦前辈一样，毫无顾忌地把分数或整数与分数之和当作数来处理。此外，他们之所以把这一限制加于数的观念，似乎更多是出于哲学的而非数学的动机，是承袭最早一批希腊思想家关于一与多的思考的结果，因为在这种思想体系中，单位一经细分，便会因这一行为本身而失去其作为单位的性质。\(^{3}\)
\item
  量的理论以公理化方式建立，而且同时对一切种类的量建立（其中可以找到对先前一些理论的暗示，那些理论看来是分别处理长度、面积、体积、时间等的）。某一类量的特征在于：它们可以比较（也就是说，假定了一个严格说来是等价关系的相等定义，以及关系 \textgreater{} 和 \textless{}）、可以加减（\(A + B\) 有定义，而且当 A \textgreater{} B 时 A - B 也有定义），并满足所谓“阿基米德”公理；这条公理从一开始就被清楚地构想为整个大厦的拱顶石（事实上，对实数的一切公理化刻画，它都是不可缺少的）；它得到阿基米德公理这个名字纯属偶然，而阿基米德本人在其《抛物线求积》的引言（{[}5 b{]}，第 II 卷，第 265 页）中就强调：这条公理已被他的前辈们使用过，它在欧多克索斯的工作中起着本质的作用，而且它的推论并不比那些未借助它而作出的面积和体积测定更缺少确定性。\(^{4}\)
\end{enumerate}

容易看出，实数的理论必然从这个公理基础中得出。应当指出，对欧多克索斯来说，给定种类的量构成一个具有单一内部合成法则（加法）的系统，但这个系统还有一个以量的比为算子的外部合成法则，而这些比被设想为构成一个交换的乘法群。设 A 与 \(A'\) 是同一种类的量，B 与 \(B'\) 亦然；A 与 \(A'\) 之比和 B 与 \(B'\) 之比被定义为相等，如果无论整数 m 与 \(m'\) 为何，\(mA < m'A'\) 都蕴含 \(mB < m'B'\)，且 \(mA > m'A'\) 都蕴含 \(mB > m'B'\)；比之间的不等式以类似的方式定义。这些比对每一种类的量都构成一个算子域，这一点等价于第四比例项存在的公理（欧几里得的行文中未明确陈述但经常使用）：给定了比 \(A/A'\) 并给定了 \(B'\)，就存在一个与 \(B'\) 同种类的 B，使得 \(B/B' = A/A'\)。这样，欧多克索斯那个极具独创性的想法使得对各种类的量分别定义的算子域得以相互等同；⁵类似地，整数集（见前文）可以等同于量的比集的一个子集，即有理比（可公度量之比）的集合；然而，由于这些比作为整数上的算子（一般）只对整数集的一个子集有定义，所以仍有必要对它们的理论单独加以展开（欧几里得第 VII 卷）。

这样构造出来的普遍算子域，对希腊数学家来说就相当于我们所说的实数集；此外同样清楚的是，凭借量的加法和量的比的乘法，他们拥有了相当于我们的实数域的东西，尽管其形式远不如我们的便利。\(^{6}\)另一方面，人们也可以自问，他们是否曾把这些集合（给定种类的量的集合，或量的比的集合）设想为我们意义上的完备集合；否则的话，就不容易明白他们为什么会（甚至不觉得有必要把它立为公理）假定第四比例项的存在；再者，某些文本看来涉及这类想法；最后，他们当然把“一条可由连续运动描画的曲线，不能从直线的一侧过渡到另一侧而不与直线相交”视为显然，这一原则他们曾在关于倍立方的研究（通过曲线的交作出 \(\sqrt[3]{2}\)）中使用过，而且它与所考虑的性质本质上等价；然而，我们掌握的文本不容许我们完全精确地了解他们在这一点上的想法。

这就是希腊数学古典时代实数理论的状况。欧多克索斯的构造尽管令人赞叹，从连贯性和严格性的观点看无可指摘，但必须承认它缺乏柔性，很不利于数值计算尤其是代数计算的发展。更何况，它的逻辑必然性只有崇尚严格、精于抽象的头脑才能领会；因此，随着希腊数学的衰落，人们看到由计算术家的传统保留下来的“朴素”观点逐渐重新露头，这很自然；例如在丢番图{[}91 a{]}那里占上风的就是这种观点，他是这一传统真正的传人，而非官方希腊科学的代表；后者虽然出于形式上的考虑复述了欧几里得的数的定义，但他用“数”这个词真正指的，是代数问题中的未知数，其解或为整数，或为分数，甚至为无理数。\(^{7}\)尽管关于数这一主题的态度上的这种变化，联系着数学史上最重要的进展之一，即代数的发展（见第 48 页），但必须明白，它本身并不构成进步，反而是退步。

我们不可能在这里追述直到中世纪末为止数的观念在印度、阿拉伯和西方数学中经历的一切变迁：占据统治地位的是数的“朴素”概念；而且，尽管这一时期欧几里得的《几何原本》被用作数学教学的基础，欧多克索斯的学说很可能普遍不曾被人理解，因为对它的需要已不再出现。欧几里得的“比”最常被冠以“数”之名；整数的计算法则被应用于它们，并由此得到正确的结果，却不去深入剖析这些方法何以成功的原因。

然而我们看到，早在 16 世纪中叶，R. 邦贝利就在其《代数学》{[}28 b{]}\(^{8}\)中就这一问题写下了一个观点，它（在假定欧几里得第 V 卷的结果已经掌握的条件下）本质上是正确的；他认识到，一旦选定了长度单位，长度与量的比之间便存在一个双射对应，于是他在长度上定义各种代数运算（当然假定单位固定），并且用长度表示数，从而得到了实数域的几何定义（这一观点的功劳通常归于笛卡尔），这样就给他的《代数学》奠定了一个坚实的几何基础。\(^{9}\)

但邦贝利的《代数学》尽管对那个时代而言先进得出奇，也没有超出根式开方和二、三、四次方程的根式求解；当然，开根式的可能性被他毫无讨论地接受了下来。西蒙·斯蒂文{[}295{]}在数的问题上也采取类似的观点：数对他而言就是标记量之量度的东西，而且他认为它本质上是“连续的”（并未明确他对这个词所赋予的含义）；即使他区分“几何数”和“算术数”，那也只是出于其定义方式的偶然，在他看来二者并无种类之别；下面就是他关于这个问题的最后的话：“因此我们断定，不存在什么荒谬的、无理的、不规则的、不可言喻的或聋的数；数之中有如此卓越与和谐，以至它们那令人赞叹的完美给了我们日夜沉思的题材”（{[}295{]}，第 10 页）。另一方面，由于第一个把十进小数这件工具变成一种计算方法，并为它们提出了一种已接近我们今天所用的记号，他清楚地认识到这些小数为每个实数提供了一个无限逼近的算法，正如他在 1594 年的那本题为“包含一条适用于一切方程的通则”的《代数附录》（该小册子唯一已知的一册于 1914 年在鲁汶被焚毁；但可参见{[}295{]}，第 88 页）中所表明的。把这样一个方程写成 \(P(x) = Q(x)\) 的形式（其中 \(P\) 是次数高于多项式 \(Q\) 的多项式，且 \(P(0) < Q(0)\)）后，我们对 \(x\) 代入 10、100、1000、\ldots{} 直到发现 \(P(x) > Q(x)\)，他说，这就确定了根的位数；然后（比如若发现根有两位）代入 10、20、\ldots{} 这就确定了十位上的数字；接着对下一位数字如法炮制，然后对相继的各位小数也是如此：“如此无限进行下去”，他说，“就会无限接近所求者”（{[}295{]}，第 88 页）。可以看出，斯蒂文（无疑是最早）对波尔查诺定理有了清晰的想法，并且认识到这一定理是数值方程系统求解的本质工具；与此同时，在那里还可以认出一种关于数连续统的直观概念，它如此清晰，以至要把它最终精确化已所剩无几。

然而，在随后的两个世纪里，正确方法的最终确立两度被两门理论的发展所延缓，我们无须在此叙述它们的历史：无穷小演算和级数理论。在它们引发的争论中，我们可以辨认出——如同在数学史的一切时期一样——两类人之间永恒的拉锯：一类是埋头向前的研究者，不惜付出一些不安全的代价，并相信日后总会有时间去巩固业已征服的领土；另一类是批判性的头脑，他们（在直觉能力和发明才干上未必输给前者）并不认为把自己的气力花在概念的精确表述和严格证成上是白费。在 17 世纪，争论的主要对象是无穷小的概念，它凭借所能达到的结果而获得事后的辩护，却似乎与阿基米德公理公开对立；我们看到那个时代最开放的头脑最终采取了一个与邦贝利相去不远的观点，其区别主要在于对古人严格方法给予了更多注意；艾萨克·巴罗（牛顿的老师，他自己也在无穷小演算的创立中起过重要作用）在其 1664-65-66 年间于剑桥讲授的《数学讲义》{[}16 a 与 b{]}中对这一观点作了精彩的阐述；他认识到，要让数的课题重新获得那句格言式的“几何确实性”，就必须回到欧多克索斯的理论，于是他极其公正地长篇为后者辩护（据他见证，这理论在他的许多同时代人看来是可理解的），以驳斥那些以晦涩甚至荒谬相责难的人。另一方面，他把数定义为表示量之比的符号，它们可借助算术运算彼此结合，由此得到了实数域，这些说法后来被牛顿在其《算术》中所承袭，而直到戴德金和康托尔，他的后继者们对之未曾改动分毫。

但恰恰是在这个时期前后，级数展开的方法被引入了，它不久就在积习难改的代数学家手中带上了纯粹形式化的性质，并把数学家的注意力从收敛问题上引开，而在实数的领域中，对级数的健全使用必然会提出收敛问题。这一方法的主要创立者牛顿仍然意识到必须考虑这些问题：而且，即使他没有把它们充分阐明，他至少已认识到，就变量的小值而言，他所引入的幂级数“在大多数情况下”至少收敛得同几何级数（其收敛性古人心知肚明）一样好（{[}233 a{]}，第 I 卷，第 3–26 页）；大约在同一时期，莱布尼茨已观察到，各项绝对值递减且趋于 0 的交错级数是收敛的；到下一个世纪，达朗贝尔于 1768 年对非收敛级数的使用表示了怀疑。但伯努利、尤其是欧拉的权威，使这类怀疑在当时只是例外。

显然，习惯于在数值计算中使用级数的数学家决不会这样忽视收敛的概念；而在这门领域中如同在许多其他领域中一样，第一个带来正确方法之回归的人，是一位从青年时代起就热爱数值计算的数学家，这并非出于偶然：C. F. 高斯几乎在孩提时代就使用过算术几何平均的算法，\(^{10}\)因而不能不为自己形成清晰的极限观念；我们看到他在一篇写于 1800 年（但直到我们时代才发表）的残稿（{[}124 a{]}，第 X \(_1\) 卷，第 390 页）中，一方面精确地定义了实数序列的上确界和下确界，另一方面定义了它的上极限和下极限；前者的存在性（对有界序列）似乎被当作显然而接受，后者则被正确地定义为当 \(n\) 趋于 +∞ 时 \(\sup_{p≥0} u_{n+p}\)、\(\inf_{p≥0} u_{n+p}\) 的极限。高斯还在其 1812 年关于超几何级数的论文（{[}124 a{]}，第 III 卷，第 139 页）中给出了关于收敛性的讨论的第一个典范，照他的说法，这讨论是“以全部的严格性进行的，为的是让那些偏爱古代几何学家严格方法的人满意”：诚然，这项讨论在该论文中只构成一个次要之点，并没有追溯到级数理论的第一原理；是柯西在其 1821 年的《分析教程》（{[}56 a{]}，(2)，第 III 卷）中最先建立了这些原理，其方式在一切点上都正确，从明确陈述并被视为显然的柯西准则出发；由于在数的定义上他拘守巴罗和牛顿的观点，因此可以说，对他而言，实数是由量的公理和柯西准则来定义的：这确实足以定义它们。

也是在这个时期，实数理论的另一个重要方面得到了最终的澄清。如前所述，两条连续曲线不能彼此穿越而不相交，这一点一直被当作几何上显然的事而接受；这条原则（经过适当的精确化）也将等价于直线是完全空间这一性质。这条原则仍然属于高斯 1799 年对达朗贝尔定理——即每个实系数多项式都有一个实的或复的根——所作“严格”证明的基础（{[}124 a{]}，第 III 卷，第 1 页）；高斯 1815 年给出的同一定理的证明（{[}124 a{]}，第 III 卷，第 31 页），与拉格朗日先前的一次尝试一样，依据的是一条类似但更简单的原则，即多项式不能变号而不变为零（我们已经见过斯蒂文使用这条原则）。1817 年，波尔查诺从柯西准则出发，对后一条原则给出了一个完全的证明，并把它作为一个关于数值变量的连续数值函数的类似定理的特例得到{[}27 c{]}。他（先于柯西）清楚地陈述了“柯西准则”，并试图用一个论证来证成它，而在缺乏实数的一切算术定义的情况下，这个论证只能是、也只可能是一个循环；但是，一旦这一点得到承认，他的工作就完全是正确的而且相当出色，因为其中不仅包含连续函数的现代定义（此处系首次给出）以及多项式连续性的证明，甚至还包含任意有界实数集的下确界存在的证明（他不谈集合，而是谈实数的性质，这二者是一回事）。另一方面，柯西在其《分析教程》（{[}56 a{]}，(2)，第 III 卷）中也定义了一元或多元数值变量的连续函数，同样证明了一元的连续函数不能变号而不变为零，而且用的论证与西蒙·斯蒂文的相同，一旦连续性有了定义，只要使用柯西准则（或者只要像柯西此时那样，假定称为“区间套”的等价原则——无穷十进小数的收敛当然只是它的一个特例），这论证自然就成为正确的。

一旦达到了这一点，数学家们剩下的工作只是通过纠正若干错误、填补若干空缺，把已获得的结果加以精确化和发展。例如，柯西一度相信，各项都是变量的连续函数的收敛级数，其和也是一个连续函数：阿贝尔在其关于级数的重要工作（{[}1{]}，第 I 卷，第 219 页；另参见第 II 卷，第 257 页及以下各处）中对这一点的纠正，最终以魏尔斯特拉斯在其讲课（未出版，但影响巨大）中对一致收敛概念的阐明而告结束（见第 205 页）。另一方面，柯西在他给出的多项式根存在性的一个证明中，曾没有充分根据地假定了连续函数最小值的存在；又是

魏尔斯特拉斯在其讲课中通过证明定义在闭有界区间上的数值变量函数的最小值存在，使这类问题得到了澄清；正是从他对把这一定理不加证成地应用于函数集（其中“狄利克雷原理”是最著名的例子）的批评出发，开始了那场思想运动，它最终（见第 143 页）以紧空间的一般定义和该定理的现代陈述而告终。

与此同时，魏尔斯特拉斯在其讲课中认识到，把实数的观念与量的理论完全分离开来在逻辑上是有意义的：这样做的确可以归结为以公理化方式定义直线上的点集（从而最终定义实数集）并假定这样一个集合存在；尽管这种做法本质上是正确的，但显然更可取的是只从有理数出发，然后通过完备化从那里导出实数。\(^{11}\)魏尔斯特拉斯、戴德金、梅雷和康托尔以不同的方法、彼此独立地做到了这一点；戴德金提出的所谓“分割”法（{[}79{]}，第 II 卷，第 315–334 页）与欧多克索斯的诸定义接近得多，而其他几种提出的方法则较接近于豪斯多夫以后用来完备化一个度量空间的方法。也正是在这个时刻，康托尔开始发展实数集的理论——戴德金早就有了关于它的最初想法{[}48{]}——并由此得到了关于直线拓扑、开集与闭集的结构、导集、完全不连通的完全集等概念的主要初等结果。

到康托尔为止，实数理论除去少数例外已具有其最终形式；我们来简要地指出它是以何种方式被推广的。除一般拓扑方面的工作（见第 143 页）以及对积分论的应用（见第 221 页）之外，这里主要涉及的是关于直线上点集及实变量的数值函数的结构与分类的研究；它们起源于博雷尔的工作{[}32 a{]}，这些工作主要着眼于测度论，但除其他结果外，最终得出了“博雷尔集”的定义：这些集属于 R 的子集中包含诸区间、且对并、可数交以及 C 运算封闭的最小族。与这些集密切相关的是所谓“贝尔”函数，也就是可以从连续函数出发、经“超穷”重复的极限运算而得到的那些函数；贝尔在一项重要工作中定义了它们，在这项工作中他完全抛弃了测度的观点，以便系统地处理这些问题的定性方面和“拓扑”方面{[}11 a{]}：正是在这个场合，他第一次定义并研究了半连续函数，而且为了刻画连续函数的极限函数，他引入了贫集（用贝尔的术语说，“第一纲”集）这一重要概念（见第 165 页）。至于继贝尔之后的大量工作——它们主要出自俄国尤其是波兰学派——我们在这里只能提请注意它们的存在（见第 166 页）。

\chapter{指数与对数.}

实数乘法群 \(R_{+}^{*}\)（\textgreater0）的历史，与一个数的幂（\textgreater0）这一观念的发展以及用来表示它们的记法密切相关。由同一个数的逐次幂组成的“几何级数”这一概念，可以追溯到埃及人和巴比伦人；希腊数学家对它也很熟悉，而且在欧几里得（《几何原本》第九卷，命题 11）那里就已经可以找到一个一般命题，它与规则 \(a^{m}a^{n}=a^{m+n}\)（指数为 \textgreater0 的整数）等价。在中世纪，法国数学家 N. 奥雷姆（十四世纪）重新发现了这条规则；分数指数（\textgreater0）的概念也首次出现在他那里，其记法已经与我们的相近，而且还伴有一些以一般形式陈述的关于它们的运算法则（例如我们现在写作 \((ab)^{1/n}=a^{1/n}b^{1/n},(a^{m})^{p/q}=(a^{mp})^{1/q}\) 的那两条法则）{[}74{]}。但奥雷姆的思想远远超前于他那个时代的数学，以致未能对其同时代人发生影响，他的论著很快就被遗忘了。一个世纪之后，N. 许凯再次陈述了欧几里得的法则；他还在自己的方程中为未知量的幂引入了一种指数记法，并且毫不犹豫地使用了指数 0 和负整数指数（\textless0）。¹这一次（尽管许凯的著作一直停留在手稿状态，似乎也未曾广泛流传），指数的“算术级数”与幂的“几何级数”之间的同构这一观念再也没有从人们的视野中消失；它被施蒂费尔推广到负指数和分数指数（{[}298{]}，fol.~35、249-250），最终导致了 对数的定义和最初数表的编制，后者由苏格兰人 J. 纳皮尔于 1614-1620 年间、瑞士人 J. 比尔吉（其著作迟至 1620 年才问世，尽管其构想可追溯到十七世纪的最初几年）各自独立地完成。在比尔吉那里，通过在处理数表时使用插值，R 与 \(R_{+}^{*}\) 之间所建立的同构的连续性是被隐含地假定的；相反，它被明确地表述于

纳皮尔的定义之中（至少在当时通行的、还相当模糊的连续性概念所允许的程度上是明确的）。\(^{2}\)

我们在这里无需详述对数在数值计算中所提供的服务；从理论观点看，它们的重要性尤其始于无穷小计算的萌芽时期，即 \(\log(1+x)\) 与 \(e^{x}\) 的级数展开以及这些函数的微分性质被发现之时（见第 172 页及以下）。至于指数函数和对数函数的定义，直到十九世纪中叶，人们都仅限于直观地假定：有可能把对一切有理数 x 定义的函数 \(a^{x}\) 连续地延拓到实数集上；而只是在实数概念被最终精确化并从有理数概念导出之后，人们才想到要对这一延拓给出严格的证成。

\chapter{n 维空间.}

我们已经有机会说过，平面和空间的解析几何的发展如何引导数学家引入 n 维空间的概念；这个概念为他们提供了一种几何语言，极其便于以简洁明了的方式表达关于任意多个变量的方程的代数定理，尤其是线性代数的一切一般结果（见第 61 页）。但是，尽管到十九世纪中叶这种语言已在众多几何学家中间通行，它仍然是纯粹约定性的，而且由于缺乏三维以上空间的“直观”表象，在这些空间中，在平面或空间里仅凭“直觉”便被允许的那种“按连续性”的论证似乎遭到了禁止。黎曼在其关于位置分析和几何学基础的研究中，第一个敢于通过与三维空间的类比按这种方式进行论证（见第 176 页）；\(^{1}\)在他之后，众多数学家开始卓有成效地使用这类论证，尤其是在多个复变量的代数函数理论中。但由于当时对空间直观的把握非常有限，人们有理由对这类思考在证明中的价值保持怀疑，只承认它们纯粹是启发性的，因为它们只是使某些定理的正确性显得颇为可能。因此，H. 庞加莱在其 1887 年关于两个复变量的函数的二重积分的留数的论文中，尽可能避免一切对四维空间直观的求助：“由于这种超几何语言至今仍令许多优秀的头脑感到反感”，他说，“我只偶尔使用它”；他为达到这一目的而使用的种种“技巧”，使他得以回到三维空间中的拓扑论证，在那里他也就不再犹豫去诉诸直观了（{[}251 a{]}，第 III 卷，第 443 页及以下）。

此外，康托尔的种种发现，尤其是确立 R 与 \(R^{n}\) 等势的那条著名定理（它似乎使维数概念本身也成了问题）\(^{2}\)表明，为了给几何学与拓扑学的论证提供坚实的基础，必须使它们完全摆脱对直观的任何求助。我们已经说过（参见第 139 页），这一需要正是一般拓扑的现代观念的起源；但甚至在这一学科创立之前，对数空间的拓扑及其最直接的推广（“n 维流形”）的严格研究就已经开始了，其所用方法主要源自拓扑学中称为“组合拓扑”或者更好地说“代数拓扑”的那个分支，关于它我们在这里无法论述。
